% TLDR: Template for experiment README.tex files. Copy this to your experiment dir and fill in.
% Reference: pass-tests-fail-proofs/experiments/01_bayes_risk_oracle_comparison/readme.tex
% Reference: veribench/experiments/38_full_metric/README.tex
%
% Every experiment README.tex should have these sections:
%   1. Objective — what question does this experiment answer?
%   2. Hypothesis — testable predictions (H1, H2, ...)
%   3. Experimental Setting — dataset, agents/models, infrastructure
%   4. Metric Definition — formal math definitions of all metrics
%   5. Results — tables with actual numbers
%   6. Why These Numbers Are Correct — sanity checks, diagnosis, interpretation
%   7. Known Limitations — honest about confounds
%   8. Verification Checklist — must-pass checks before reporting
%   9. Reproduction — exact commands to re-run
%  10. Experimental Details — model IDs, dates, file paths, W&B links

\documentclass[11pt]{article}
\usepackage[margin=1in]{geometry}
\usepackage{amsmath,amssymb,amsthm}
\usepackage{booktabs}
\usepackage{hyperref}
\usepackage{enumitem}
\usepackage{xcolor}

\title{Experiment NN: [Short Descriptive Title]}
\author{Brando Miranda \\ Stanford University \\ \texttt{brando9@stanford.edu}}
\date{Started: YYYY-MM-DD \quad|\quad Status: [In Progress / Complete]}

\begin{document}
\maketitle

%% ============================================================
\section{Objective}
\label{sec:objective}
%% ============================================================

% 1-2 paragraphs. What question does this experiment answer?
% What is the independent variable? What is the dependent variable?

TODO

%% ============================================================
\section{Hypothesis}
\label{sec:hypothesis}
%% ============================================================

% Testable predictions. Number them H1, H2, etc.
% Each should be falsifiable by the experiment's results.

\begin{quote}
\textbf{H1:} TODO --- [prediction about the outcome]

\textbf{H2:} TODO --- [second prediction]
\end{quote}

%% ============================================================
\section{Experimental Setting}
\label{sec:setting}
%% ============================================================

\subsection{Dataset}

% Table of dataset splits with counts and descriptions.

\begin{center}
\begin{tabular}{lrl}
\toprule
Split & $n$ & Description \\
\midrule
TODO & 0 & TODO \\
\midrule
\textbf{Total} & \textbf{0} & \\
\bottomrule
\end{tabular}
\end{center}

\subsection{Models / Agents Under Evaluation}

% Table listing all agents/models, their type, and status.

\begin{table}[h]
\centering
\begin{tabular}{@{}lll@{}}
\toprule
Agent / Model & Type & Status \\
\midrule
TODO & TODO & TODO \\
\bottomrule
\end{tabular}
\caption{Agents under evaluation.}
\end{table}

\subsection{Infrastructure}

% Docker, Harbor, cluster, sandboxing, anti-cheat — whatever applies.

TODO

%% ============================================================
\section{Metric Definition}
\label{sec:metric}
%% ============================================================

\subsection{Definitions}

% Formal mathematical definitions of every metric used.
% Include: symbol, full name, formula, range, edge cases.

TODO

\subsection{Composite Score (if applicable)}

% How individual factors combine. Justify the aggregation method.

TODO

%% ============================================================
\section{Results}
\label{sec:results}
%% ============================================================

% Tables with actual numbers. Include raw data + any adjustments.
% Follow the pass-tests-fail-proofs pattern:
%   1. Raw results table
%   2. "Why these numbers are correct" subsection with diagnosis
%   3. Adjusted results if applicable
%   4. Per-split breakdown

\subsection{Overall Results}

\begin{table}[h]
\centering
\begin{tabular}{@{}lcc@{}}
\toprule
Agent & Metric 1 & Metric 2 \\
\midrule
TODO & 0.000 & 0.000 \\
\bottomrule
\end{tabular}
\caption{Overall results. Date: YYYY-MM-DD. Model: \texttt{model-id}.}
\end{table}

\subsection{Why These Numbers Are Correct}

% For EACH major result, explain WHY you believe it:
%   - Sanity checks that passed (oracle, ceiling, floor)
%   - Cross-validation with prior experiments
%   - Diagnosis of surprising results (e.g., FP analysis)
%   - Consistency checks (factors that should be identical are identical)

TODO

\subsection{Per-Split Analysis (if applicable)}

TODO

%% ============================================================
\section{Known Limitations}
\label{sec:limitations}
%% ============================================================

% Honest about confounds, approximations, and threats to validity.

\begin{itemize}
  \item TODO
\end{itemize}

%% ============================================================
\section{Verification Checklist}
\label{sec:verification}
%% ============================================================

% Must-pass checks before reporting results as final.
% Use checkboxes. Agent must confirm each before committing.

Before reporting results as final:

\begin{enumerate}
  \item[$\square$] TODO --- sanity check 1
  \item[$\square$] TODO --- sanity check 2
  \item[$\square$] Result files match (no transcription errors)
\end{enumerate}

%% ============================================================
\section{Reproduction}
\label{sec:reproduction}
%% ============================================================

% Exact bash commands to reproduce from scratch.

\begin{verbatim}
# Step 1: ...
TODO

# Step 2: ...
TODO
\end{verbatim}

%% ============================================================
\section{Experimental Details}
\label{sec:details}
%% ============================================================

\begin{itemize}
  \item \textbf{Model:} TODO (\texttt{model-id})
  \item \textbf{Date:} YYYY-MM-DD
  \item \textbf{Result files:}
    \begin{itemize}
      \item \texttt{results/TODO.json}
    \end{itemize}
  \item \textbf{W\&B:} \url{https://wandb.ai/brando-su/TODO}
\end{itemize}

\end{document}
